\documentclass[letterpaper,12pt]{article}
%\usepackage{harvard}
%\usepackage{subfigure}
%\usepackage[active]{srcltx}
\usepackage{graphicx}
\usepackage{pdflscape}
\usepackage{amsmath}
\usepackage{amsthm}
\usepackage{lscape}
\usepackage{enumerate}
\usepackage{float}
\usepackage{grffile}
\usepackage{relsize}
\usepackage{tablefootnote}
\usepackage{footnote}
\usepackage{sectsty}
\usepackage{xcolor}
\usepackage{caption}
\usepackage{color}
\usepackage{natbib}
\usepackage{multirow}
\usepackage{bbm}
\usepackage[T1]{fontenc}
\usepackage[multiple]{footmisc}
\usepackage{appendix}
\usepackage{times}
\usepackage{relsize}
\usepackage{caption}
\usepackage{newtxtext,newtxmath}
\usepackage{enumitem}
\usepackage[breaklinks,colorlinks]{hyperref}

%%%%%%%%%%%%%%%%%%%%%%%%%%%%%%%%%%%%%%%%%%%%%%%%%%%%%%%%%%%%
%%%   Title Page Information
%%%%%%%%%%%%%%%%%%%%%%%%%%%%%%%%%%%%%%%%%%%%%%%%%%%%%%%%%%%%
\title{ECON311: Critical Review}
\author{Leandro Freylejer}
\date{Due date in class on March 31st, 2026}

%%%%%%%%%%%%%%%%%%%%%%%%%%%%%%%%%%%%%%%%%%%%%%%%%%%%%%%%%%%%
%%%   Double and Single Space Commands
%%%%%%%%%%%%%%%%%%%%%%%%%%%%%%%%%%%%%%%%%%%%%%%%%%%%%%%%%%%%
\newlength{\defbaselineskip}
\setlength{\defbaselineskip}{\baselineskip}
\newcommand{\setlinespacing}[1]{\setlength{\baselineskip}{#1 \defbaselineskip}}
\newcommand{\doublespacing}{\setlength{\baselineskip}{1.3 \defbaselineskip}}
\newcommand{\singlespacing}{\setlength{\baselineskip}{1\defbaselineskip}}
\renewcommand{\thesubsection}{\arabic{subsection}}
\newtheoremstyle{dotlessS}{}{}{\itshape}{}{\color{dg}\bfseries}{}{ }{}
\theoremstyle{dotlessS}
\newtheorem{theorem}{Theorem}
\newtheorem{lemma}{Lemma}
\newtheorem{proposition}{Proposition}
\newtheorem{assumption}[theorem]{Assumption}
\newtheorem{corollary}{Corollary}
\newtheorem{definition}{Definition}
\newenvironment{example}[1][Example]{\begin{trivlist}
\item[\hskip \labelsep {\bfseries #1}]}{\end{trivlist}}
\newenvironment{remark}[1][Remark]{\begin{trivlist}
\item[\hskip \labelsep {\bfseries #1}]}{\end{trivlist}}

\newcommand\Pitem{%
  \addtocounter{enumi}{-1}%
  \renewcommand\theenumi{\arabic{\enumi}'}%
  \item%
  \renewcommand\theenumi{\arabic{enumi}}%
}

%%%%%%%%%%%%%%%%%%%%%%%%%%%%%%%%%%%%%%%%%%%%%%%%%%%%%%%%%%%%
%%%   Set Margins Here
%%%%%%%%%%%%%%%%%%%%%%%%%%%%%%%%%%%%%%%%%%%%%%%%%%%%%%%%%
\textwidth=6.5in
\textheight=9in
\oddsidemargin=0.10in
\evensidemargin=0.10in
\topmargin=-0.5in
%\parskip=3pt

%%%%%%%%%%%%%%%%%%%%%%%%%%%%%%%%%%%%%%%%%%%%%%%%%%%%%%%%%%%%%
%%%   Actual Document Begins Here
%%%%%%%%%%%%%%%%%%%%%%%%%%%%%%%%%%%%%%%%%%%%%%%%%%%%%%%%%%%%%
\begin{document}
\pagenumbering{arabic}
\maketitle
\vspace{-.3in}
\section*{General Task}
For this writing assignment you are asked to critically analyze the article titled ``The COVID crisis and productivity growth''. There are two crucial components to writing a critical review. First, you must provide a clear summary of what you consider the main points of the article. Second, you must provide a critical analysis of these points. Your critical analysis should use concepts that you have learned in this class, other economics classes, and your own additional research on the topic. The final version of your review should be between \textbf{five} and \textbf{six} double-spaced pages in length.  \\
\\
Critical Reviews are more than merely paraphrasing the author's main points. It involves a careful examination of an article within the broader context provided by your knowledge of the topic being analyzed. Therefore, it is very important that you consult additional references discussing the main ideas of the article. You will be graded on your general understanding of the article, the thoroughness of your analysis, and the clarity of your argument.\\
\\
A useful, but not necessary, structure for your analysis will start with an introduction that provides some context and discusses what you consider the key points of the article. In the middle part of the critical review, you should analyze the arguments of the article using the knowledge that you have acquired in your own research. Finally, the conclusion should briefly summarize your critical review and include a concluding sentence on the quality of the article. Any reference you cite should be identified at the end of the paper. As a side note, though the review is called the Critical Review, your discussion should highlight both the strengths and weaknesses of the article.               

\section*{Format of the Assignment}
\begin{itemize}
\item Your assignment should be typed, double spaced and have a font size of 12 points. 
\item The assignment should be either a Word or pdf document.
\item You must thoroughly cite your sources using the ``APA System''. You can find detailed information about employing the ``APA System'' to cite your references at \url{https://www.unbc.ca/academic-success-centre/online-writing-resources}
\item Please do not either intentionally or inadvertently plagiarize. The Academic Success Centre at UNBC is a great resource to learn more about plagiarism and writing in general. You can find many handouts with information in the link provided above. 
\item The use of Generative AI is strictly prohibited (except as a search engine).    
\end{itemize}

\section*{Grading Rubric}
The Critical Review will be graded out of 50 points as follows

\begin{itemize}
\item \textbf{Introduction [out of 5 points]}
\begin{itemize}
\item \textbf{Excellent [4-5]}
\begin{itemize}
\item Provides a clear context to the article
\item Presents a clear and accurate representation of the author's main arguments
\item Sets up the main thesis of the review article
\end{itemize}
\item \textbf{Good [3-4]}
\begin{itemize}
\item Provides some context to the article
\item Presents an accurate representation of the author's main arguments
\item The main thesis of the review article is not entirely clear
\end{itemize}
\item \textbf{Competent [2-3]}
\begin{itemize}
\item Provides no context to the article
\item Presents an inadequate representation of the author's main arguments
\item There is no main thesis of the review article 
\end{itemize}
\item \textbf{Problematic [0-2]}
\begin{itemize}
\item Provides no context to the article
\item The author's main arguments are inaccurately represented
\item There is no main thesis of the review article 
\end{itemize}
\end{itemize}


\item \textbf{Main Analysis [out of 30 points]}
\begin{itemize}
\item \textbf{Excellent [25-30]}
\begin{itemize}
\item Provides a clear and detailed critical analysis of the author's main arguments using additional material
\item Shows independent critical thinking rather than paraphrasing additional material    
\item Criticism concentrates on important economic points and not minor aspects
\item Shows depth of knowledge of the subject acquired through additional research
\end{itemize}
\item \textbf{Good [20-25]}
\begin{itemize}
\item Provides a clear critical analysis of the author's main arguments using additional material
\item Shows some independent critical thinking rather than paraphrasing additional material    
\item Criticism concentrates on somewhat economic points and not minor aspects
\item Shows some knowledge of the subject acquired through additional research
\end{itemize}
\item \textbf{Competent [15-20]}
\begin{itemize}
\item Provides an unclear critical analysis of the author's main arguments using additional material
\item Mostly paraphrases additional material    
\item Criticism concentrates on minor aspects
\item Shows very little knowledge of the subject acquired through additional research
\end{itemize}
\item \textbf{Problematic [0-15]}
\begin{itemize}
\item Provides no critical analysis of the author's main arguments using additional material
\item Mostly paraphrases additional material    
\item Shows no additional independent research
\end{itemize}
\end{itemize}

\item \textbf{Conclusion [out of 5 points]}
\begin{itemize}
\item \textbf{Excellent [4-5]}
\begin{itemize}
\item Summarizes the analysis briefly and accurately
\end{itemize}
\item \textbf{Good [3-4]}
\begin{itemize}
\item Attempts to summarize the analysis briefly and accurately
\end{itemize}
\item \textbf{Competent [2-3]}
\begin{itemize}
\item Summarizes the analysis somewhat accurately
\end{itemize}
\item \textbf{Problematic [0-2]}
\begin{itemize}
\item The conclusion does not summarizes the analysis.
\end{itemize}
\end{itemize}



\item \textbf{Exposition [out of 10 points]}
\begin{itemize}
\item \textbf{Excellent [8-10]}
\begin{itemize}
\item Excellent organization of the discussion
\item Well-written paragraphs with smooth transitions between sentences and across paragraphs
\item No unnecessary repetition of arguments
\item Clear explanations
\item Sources are referenced and cited properly
\end{itemize}
\item \textbf{Good [7-8]}
\begin{itemize}
\item Good organization of the discussion
\item Well-written paragraphs with smooth transitions between sentences and across paragraphs
\item Some unnecessary repetition of arguments
\item Explanation not always clear
\item Sources are referenced and cited properly
\end{itemize}
\item \textbf{Competent [5-7]}
\begin{itemize}
\item Poor organization of the discussion
\item Paragraphs are not always clear 
\item Some unnecessary repetition of arguments 
\item Disorganize explanation
\item Sources are sometimes referenced and cited properly 
\end{itemize}
\item \textbf{Problematic [0-5]}
\begin{itemize}
\item Poor organization of the discussion
\item Paragraphs are not clear
\item Paraphrases without providing context 
\item Disorganize explanation
\item Sources are incorrectly referenced and cited
\end{itemize}
\end{itemize}
\end{itemize}

\end{document}










